% Faithful translation baseline; Chinese source PDF pp. 35--40.
\section{One-Dimensional Linear Harmonic Oscillator}
For an arbitrary potential, expand it in a Taylor series near a minimum $x_0$:
\begin{equation}
V(x)=V(x_0)+V'(x_0)(x-x_0)+\frac12V''(x_0)(x-x_0)^2+\cdots.
\end{equation}
\begin{center}
\begin{tikzpicture}[>=Latex,scale=.85]
\draw[->](0,0)--(4.9,0) node[right]{$x$};\draw[->](0,0)--(0,2.5)node[above]{$V(x)$};
\draw[thick,smooth] plot coordinates {(.6,2.0)(.9,1.35)(1.2,1.0)(1.5,1.42)(1.95,1.62)(2.7,1.8)(3.4,1.87)(4.4,1.68)};
\draw[dashed](1.2,0)--(1.2,1.0);\node[below]at(1.2,0){$x_0$};
\end{tikzpicture}
\end{center}
The constant term $V(x_0)$ can be incorporated into the energy. At the potential minimum $V'(x_0)=0$. Omitting higher-order terms,
\begin{equation}
V(x)\simeq\frac12V''(x_0)(x-x_0)^2.
\end{equation}
This approximates a harmonic-oscillator potential; therefore the harmonic oscillator has general significance.

Classically,
\begin{equation}
H=\frac{p^2}{2m}+\frac12m\omega^2x^2.
\end{equation}
Quantum mechanically,
\begin{equation}
\op H=\frac{\op p^2}{2m}+\frac12m\omega^2\op x^2,\qquad
[\op x,\op p]=\ii\hbar.
\end{equation}
Because
\begin{equation}
\op H=\frac{m\omega^2}{2}
\left(\op x-\frac{\ii\op p}{m\omega}\right)
\left(\op x+\frac{\ii\op p}{m\omega}\right)
+\frac12\hbar\omega,
\end{equation}
define new operators
\begin{equation}
\op a=\sqrt{\frac{m\omega}{2\hbar}}
\left(\op x+\frac{\ii\op p}{m\omega}\right),\qquad
\op a^{\dagger}=\sqrt{\frac{m\omega}{2\hbar}}
\left(\op x-\frac{\ii\op p}{m\omega}\right).
\end{equation}
Then
\begin{equation}
\boxed{\op H=\left(\op a^{\dagger}\op a+\frac12\right)\hbar\omega,\qquad
[\op a,\op a^{\dagger}]=\mathbf1.}
\end{equation}
Clearly, $\op a$ is not Hermitian, $\op a^{\dagger}\ne\op a$. However, $\op a^{\dagger}\op a$ is Hermitian:
\begin{equation}
(\op a^{\dagger}\op a)^{\dagger}=\op a^{\dagger}\op a.
\end{equation}
Since $\op H$ differs from $\op a^{\dagger}\op a$ only by a constant, they have common eigenstates. Let
\begin{equation}
\op a^{\dagger}\op a\ket n=n\ket n.
\end{equation}
Then
\begin{equation}
\op H\ket n=\left(n+\frac12\right)\hbar\omega\ket n,\qquad
E_n=\left(n+\frac12\right)\hbar\omega.
\end{equation}
\textbf{Questions:} What are the eigenvalues $n$? What are the coordinate-representation eigenstates $\braket x n$?

\subsection*{1) Nonnegativity of the eigenvalues}
Let
\begin{equation}
\op a\ket n=\ket b.
\end{equation}
Then
\begin{equation}
\bra n\op a^{\dagger}=\bra b,\qquad
\mel n{\op a^{\dagger}\op a}n=\braket b b,\quad
n\braket n n=\braket b b.
\end{equation}
Because $\braket b b\ge0$ and $\braket n n\ge0$, we have $n\ge0$.

\subsection*{2) Raising and lowering operators}
From $[\op a,\op a^{\dagger}]=1$,
\begin{equation}
\begin{split}
(\op a^{\dagger}\op a)\op a\ket n
&=\op a(\op a\op a^{\dagger}-1)\ket n
=\op a(\op a^{\dagger}\op a-1)\ket n\\
&=(n-1)\op a\ket n,
\end{split}
\end{equation}
\begin{equation}
\begin{split}
(\op a^{\dagger}\op a)\op a^{\dagger}\ket n
&=\op a^{\dagger}(\op a\op a^{\dagger})\ket n
=\op a^{\dagger}(\op a^{\dagger}\op a+1)\ket n\\
&=(n+1)\op a^{\dagger}\ket n.
\end{split}
\end{equation}
If $\ket n$ is an eigenstate of $\op a^{\dagger}\op a$, so are $\op a\ket n$ and $\op a^{\dagger}\ket n$, with the following relationships:
\begin{equation}
\begin{array}{c|c}
\text{Eigenstate}&\text{Eigenvalue of }\op a^{\dagger}\op a\\\hline
(\op a^{\dagger})^2\ket n&n+2\\
\op a^{\dagger}\ket n&n+1\\
\ket n&n\\
\op a\ket n&n-1\\
\op a^2\ket n&n-2
\end{array}
\end{equation}
Hence $\op a$ is called the lowering (annihilation) operator and $\op a^{\dagger}$ the raising (creation) operator. Together with $n\ge0$, the possible eigenvalues have the form
\begin{equation}
n=n_0,n_0+1,n_0+2,\ldots,\qquad n_0\ge0.
\end{equation}

\subsection*{3) Lowest eigenvalue}
For the lowest eigenvalue $n_0$, necessarily
\begin{equation}
\op a\ket{n_0}=0.
\end{equation}
Otherwise $\op a\ket{n_0}$ would have eigenvalue $n_0-1$, contradicting the assumption that $n_0$ is the smallest eigenvalue. From
\begin{equation}
\op a^{\dagger}\op a\ket{n_0}=n_0\ket{n_0}=0
\end{equation}
we obtain $n_0=0$. Therefore
\begin{equation}
\boxed{E_n=\left(n+\frac12\right)\hbar\omega,\quad n=0,1,2,\ldots.}
\end{equation}
Notice that so far we have used only the commutation relation, without introducing any specific representation.

\subsection*{4) Matrix elements of creation and annihilation operators}
Since $\op a\ket n$ has eigenvalue $n-1$, and in the absence of degeneracy (one-dimensional bound states are nondegenerate),
\begin{equation}
\op a\ket n=a_n\ket{n-1}.
\end{equation}
Similarly,
\begin{equation}
\op a^{\dagger}\ket n=b_n\ket{n+1}.
\end{equation}
Because
\begin{equation}
\bra n\op a^{\dagger}=a_n^*\bra{n-1},\qquad
\bra n\op a=b_n^*\bra{n+1},
\end{equation}
we obtain
\begin{equation}
\mel n{\op a^{\dagger}\op a}n=|a_n|^2\braket{n-1}{n-1},\quad n=|a_n|^2,\quad |a_n|=\sqrt n,
\end{equation}
\begin{equation}
\mel n{\op a\op a^{\dagger}}n=|b_n|^2\braket{n+1}{n+1},\quad
n+1=|b_n|^2,\quad |b_n|=\sqrt{n+1}.
\end{equation}
Choose the phases so that $a_n=\sqrt n$ and $b_n=\sqrt{n+1}$. Thus
\begin{equation}
\boxed{\op a\ket n=\sqrt n\ket{n-1},\qquad
\op a^{\dagger}\ket n=\sqrt{n+1}\ket{n+1}.}
\end{equation}

\subsection*{5) Representation common to $\op H$ and $\op a^{\dagger}\op a$}
\begin{equation}
(a^{\dagger}a)_{mn}=\mel m{\op a^{\dagger}\op a}n
=n\braket m n=n\delta_{mn},
\end{equation}
\begin{equation}
H_{mn}=\mel m{\op H}n=\left(n+\frac12\right)\hbar\omega\delta_{mn}.
\end{equation}
Both are diagonal matrices, as expected in their own representation. Meanwhile,
\begin{equation}
a_{mn}=\mel m{\op a}n=\sqrt n\braket m{n-1}=\sqrt n\delta_{m,n-1},
\end{equation}
\begin{equation}
(a^{\dagger})_{mn}=\mel m{\op a^{\dagger}}n
=\sqrt{n+1}\braket m{n+1}=\sqrt{n+1}\delta_{m,n+1}.
\end{equation}
Since
\begin{equation}
\op x=\sqrt{\frac\hbar{2m\omega}}(\op a^{\dagger}+\op a),\qquad
\op p=\ii\sqrt{\frac{m\hbar\omega}{2}}(\op a^{\dagger}-\op a),
\end{equation}
it follows that
\begin{equation}
x_{mn}=\sqrt{\frac\hbar{2m\omega}}
\bigl(\sqrt n\delta_{m,n-1}+\sqrt{n+1}\delta_{m,n+1}\bigr),
\end{equation}
\begin{equation}
p_{mn}=\ii\sqrt{\frac{m\hbar\omega}{2}}
\bigl(-\sqrt n\delta_{m,n-1}+\sqrt{n+1}\delta_{m,n+1}\bigr).
\end{equation}
They are nondiagonal matrices. Note that $\mel n{\op x}n=\mel n{\op p}n=0$.

\subsection*{6) Coordinate representation}
For the ground state $\ket0$,
\begin{equation}
\op a\ket0=0,\qquad
\left(\op x+\frac{\ii}{m\omega}\op p\right)\ket0=0.
\end{equation}
Multiplying from the left by $\bra x$ and inserting the completeness relation,
\begin{equation}
\int\dd x'\,\bra x\left(\op x+\frac{\ii}{m\omega}\op p\right)\ket{x'}\braket{x'}0=0.
\end{equation}
Using
\begin{equation}
\mel x{\op x}{x'}=x\delta(x-x'),\qquad
\mel x{\op p}{x'}=-\ii\hbar\frac{\dd}{\dd x}\delta(x-x'),\qquad
\braket{x'}0=\psi_0(x'),
\end{equation}
we obtain
\begin{equation}
\left(x+\frac\hbar{m\omega}\frac{\dd}{\dd x}\right)\psi_0(x)=0.
\end{equation}
The solution of this first-order differential equation is
\begin{equation}
\psi_0(x)=C\exp\left(-\frac{m\omega x^2}{2\hbar}\right).
\end{equation}
With the normalization condition
\begin{equation}
\int_{-\infty}^{\infty}|\psi_0(x)|^2\dd x=1,
\end{equation}
we have
\begin{equation}
\boxed{\psi_0(x)=\left(\frac{m\omega}{\pi\hbar}\right)^{1/4}
\exp\left(-\frac{m\omega x^2}{2\hbar}\right).}
\end{equation}
For an excited state,
\begin{equation}
\begin{split}
\psi_n(x)=\braket x n
&=\frac1{\sqrt n}\bra x\op a^{\dagger}\ket{n-1}
=\frac1{\sqrt{n(n-1)}}\bra x(\op a^{\dagger})^2\ket{n-2}\\
&=\frac1{\sqrt{n!}}\bra x(\op a^{\dagger})^n\ket0
=\frac1{\sqrt{n!}}\left(\frac{m\omega}{2\hbar}\right)^{n/2}
\bra x\left(\op x-\frac{\ii}{m\omega}\op p\right)^n\ket0.
\end{split}
\end{equation}
Insert $n$ position completeness relations:
\begin{equation}
\begin{split}
\psi_n(x)=\frac1{\sqrt{n!}}\left(\frac{m\omega}{2\hbar}\right)^{n/2}
\int\dd x_1\cdots\dd x_n\;&\mel x{\op x-\frac{\ii}{m\omega}\op p}{x_1}
\mel{x_1}{\op x-\frac{\ii}{m\omega}\op p}{x_2}\\[-2pt]
&\cdots\mel{x_{n-1}}{\op x-\frac{\ii}{m\omega}\op p}{x_n}\braket{x_n}0.
\end{split}
\end{equation}
Using the position and momentum matrix elements gives
\begin{equation}
\boxed{\psi_n(x)=\frac1{\sqrt{n!}}\left(\frac{m\omega}{2\hbar}\right)^{n/2}
\left(x-\frac\hbar{m\omega}\frac{\dd}{\dd x}\right)^n\psi_0(x).}
\end{equation}
Thus the one-dimensional harmonic-oscillator problem has been solved completely.

\subsection*{7) Physical meaning of $\op a^{\dagger}\op a$}
The harmonic-oscillator energies are
\begin{equation}
E_n=\left(n+\frac12\right)\hbar\omega,\qquad n=0,1,2,3,\ldots.
\end{equation}
The ground-state energy $E_0=\hbar\omega/2>0$ is the zero-point energy. An excited state involves quanta of energy $\varepsilon=\hbar\omega$. When the oscillator is in the $n$-th excited state, $n$ quanta have been excited.

$n$: quantum number, particle number.

$\op N=\op a^{\dagger}\op a$ is the particle-number operator; $\op a$ is a particle-annihilation operator and $\op a^{\dagger}$ a particle-creation operator.

The common representation of $\op H$ and $\op N$ is the particle-number representation.

To summarize the algebraic solution: bound states $\to$ discrete spectrum $\to$ identify lowering and raising operators for the discrete quantum numbers. The particle-number representation is the foundation of the second-quantization scheme.

\subsection*{8) Casimir effect}
Does the zero-point energy have observable effects? Consider the vacuum energy of the quantized electromagnetic field.

According to Planck's studies of blackbody radiation, the electromagnetic field consists of harmonic oscillations of different frequencies $\omega_k$ (energy $\hbar\omega_k$). Here $k$ is the wave number, $\omega_k=ck$, and $n$ is the number of photons of energy $\hbar\omega_k$. The total electromagnetic-field energy is written
\begin{equation}
E=2\int\frac{\dd^3k}{(2\pi)^3}\sum_{n=0}^{\infty}
\left(n+\frac12\right)\hbar\omega_k.
\end{equation}
Boundary conditions change the distribution of frequencies, hence the distribution of wave numbers $k$.

The vacuum is the unexcited state, $n=0$, with vacuum energy
\begin{equation}
E=\int\frac{\dd^3k}{(2\pi)^3}\hbar\omega_k,
\end{equation}
which is infinite. In ordinary calculations it is subtracted as the background to excitations. The question is whether the vacuum energy itself has observable effects.

Place two parallel plates in the electromagnetic radiation field, each of area $L^2$ with $L$ sufficiently large and separation $a\ll L$. Boundary conditions restrict the component of wave number perpendicular to the plates to $k_z=n\pi/a$; in the transverse direction there is no significant restriction because $L$ is large. The vacuum energy between the plates, in volume $L^2a$, is
\begin{equation}
E(a)=\hbar c\int\frac{\dd^2 k_\perp}{(2\pi/L)^2}
\sum_{n=0}^{\infty}\sqrt{k_\perp^2+\left(\frac{n\pi}{a}\right)^2}.
\end{equation}
Here $\dd^2k_\perp/(2\pi/L)^2$ is the number of transverse modes in the element $\dd^2k_\perp$ near $\vect k_\perp$.

The change in vacuum energy associated with inserting the plates into this volume is
\begin{equation}
\begin{split}
\Delta E(a)=E(a)-E_0(a)
=\hbar c\int\frac{\dd^2 k_\perp}{(2\pi/L)^2}
\bigg[&\sum_{n=0}^{\infty}\sqrt{k_\perp^2+\left(\frac{n\pi}{a}\right)^2}\\[-2pt]
&-\int_0^\infty\dd n\,\sqrt{k_\perp^2+\left(\frac{n\pi}{a}\right)^2}\bigg].
\end{split}
\end{equation}
This change produces a force per unit area between the plates,
\begin{equation}
F(a)=-\frac1{L^2}\frac{\partial\Delta E(a)}{\partial a}.
\end{equation}
Both $E(a)$ and $E_0(a)$ diverge (vacuum divergence), but $\Delta E(a)$ is finite. After calculation (see standard quantum-mechanics textbooks, for example those by Su Ru-Keng and Zhang Yong-De),
\begin{equation}
\boxed{F(a)=-\frac{\hbar c\pi^2}{240a^4}.}
\end{equation}
The attractive force caused by the change (fluctuation) in vacuum energy when the plates are inserted, and its dependence on $a$, were experimentally observed in 1958. This is an observable effect of the vacuum energy.

\subsection*{9) Time evolution of observables in the Heisenberg picture}
Omitting the H subscript, the Heisenberg equations of motion give
\begin{equation}
\frac{\dd\op a}{\dd t}=\frac1{\ii\hbar}[\op a,\op H]=-\ii\omega\op a,\qquad
\frac{\dd\op a^{\dagger}}{\dd t}=\frac1{\ii\hbar}[\op a^{\dagger},\op H]
=\ii\omega\op a^{\dagger}.
\end{equation}
The solutions are
\begin{equation}
\op a(t)=\op a(0)e^{-\ii\omega t},\qquad
\op a^{\dagger}(t)=\op a^{\dagger}(0)e^{\ii\omega t}.
\end{equation}
Alternatively,
\begin{equation}
\begin{split}
\op x(t)&=\sqrt{\frac\hbar{2m\omega}}\bigl(\op a(t)+\op a^{\dagger}(t)\bigr)
=\op x(0)\cos\omega t+\frac{\op p(0)}{m\omega}\sin\omega t,\\
\op p(t)&=\ii\sqrt{\frac{m\hbar\omega}{2}}\bigl(\op a^{\dagger}(t)-\op a(t)\bigr)
=\op p(0)\cos\omega t-m\omega\op x(0)\sin\omega t.
\end{split}
\end{equation}
Although $\op x(t)$ and $\op p(t)$ oscillate, their expectation values in an energy eigenstate are
\begin{equation}
\mel n{\op x(t)}n
=\mel n{\op x(0)}n\cos\omega t+
\frac1{m\omega}\mel n{\op p(0)}n\sin\omega t=0,
\end{equation}
\begin{equation}
\mel n{\op p(t)}n
=\mel n{\op p(0)}n\cos\omega t-m\omega\mel n{\op x(0)}n\sin\omega t=0
\end{equation}
and thus show no oscillation.